\subsection*{Methods insert: additional software validation}
\paragraph{Comparison with interpolation-based resampling baselines.}
We generated deterministic synthetic anisotropic three-dimensional phantoms containing a binary ROI, an in-plane intensity gradient, through-plane alternating bands, local heterogeneous texture, and fixed-seed noise. Three through-plane spacing ratios were evaluated: 2:1:1, 3:1:1, and 5:1:1 in array order $(z,y,x)$. For each phantom, radiomic texture features were extracted using native original PyRadiomics without resampling, isotropic resampling with nearest-neighbor, linear, or B-spline image interpolation (mask resampling always nearest-neighbor), and the modified voxel-spacing-aware implementation using \texttt{weightingNorm=voxelSpacing}. Feature families were restricted to GLCM, GLRLM, GLDM, GLSZM, and NGTDM.

\paragraph{Computational profiling.}
Runtime and approximate memory use were measured for voxel-spacing-aware extraction across isotropic and anisotropic spacing settings and three volume sizes: $16\times64\times64$, $32\times96\times96$, and $48\times128\times128$. The corresponding native ROI sizes were 12,126, 50,328, and 139,028 voxels. The finite-volume expansion factor was defined as $R=q_z q_y q_x$, where $q_i=\mathrm{round}(\Delta_i/h)$ and $h=\min(\Delta_z,\Delta_y,\Delta_x)$. The represented spacing and relative error were recorded as
\[
\Delta_i^*=q_i h,\qquad e_i=\frac{\Delta_i^*-\Delta_i}{\Delta_i}.
\]
The expanded lattice size was estimated as $N_{\mathrm{expanded}}=(N_zq_z)(N_yq_y)(N_xq_x)$.

\paragraph{Finite-volume rounding and binning sensitivity.}
To evaluate non-integer spacing ratios, we tested through-plane spacings of 2.0, 2.5, 3.0, 3.5, 5.0, and 5.5 mm with 1 mm in-plane spacing. We recorded the represented spacing implied by integer repeat factors and compared neighboring feature sets. Binning sensitivity was assessed at binWidth values of 5, 10, 25, and 50 for native extraction, linear resampling, and voxel-spacing-aware extraction.

\subsection*{Results insert}
The interpolation-baseline experiment showed that voxel-spacing-aware extraction produced feature values that were not identical to either native index-space extraction or conventional isotropic resampling. Relative to native extraction, nearest-neighbor resampling changed 192/225 features (85.3\%), linear and B-spline resampling changed 225/225 features (100\%), and voxel-spacing-aware extraction changed 192/225 features (85.3\%). The median of family-level median relative differences was 0.333 for nearest-neighbor resampling, 0.304 for linear resampling, 0.302 for B-spline resampling, and 0.307 for voxel-spacing-aware extraction. Minimum method-level Spearman correlations were 0.900 for nearest-neighbor and voxel-spacing-aware extraction and 0.800 for linear and B-spline resampling. The per-feature supplementary CSV contains one row per feature with absolute and relative differences for native, nearest-neighbor, linear, B-spline, and voxel-spacing-aware extraction.

Computational profiling demonstrated that runtime and memory increased with the finite-volume expansion factor and ROI size. The most demanding completed case used a $48\times128\times128$ input with spacing $0.7{:}0.7{:}5.0$, producing $q=(1,1,7)$, an expanded lattice of $48\times128\times896$, 973,196 ROI voxels after expansion, a runtime of 1.74 s, and a peak memory estimate of 90.55 MB. In practical use, an expansion factor greater than 8 should trigger a runtime warning, and an expansion factor greater than 16 should require explicit user confirmation or a hard guard depending on deployment context.

Finite-volume rounding introduced a quantifiable approximation for non-integer spacing ratios because the current implementation maps physical spacing ratios to integer repeat factors. The examples $2.5{:}1{:}1\rightarrow2{:}1{:}1$ and $3.5{:}1{:}1\rightarrow4{:}1{:}1$ gave relative represented-spacing errors of $-20.0\%$ and $+14.3\%$, respectively. Across adjacent spacing-ratio comparisons, the largest family-level median relative feature difference was 0.167 for GLSZM and the lowest Spearman correlation was 0.900 for NGTDM. This behavior is transparent and measurable; future implementations may use rational approximations, direct physical traversal, or streaming finite-volume construction without explicit replication.

Binning sensitivity analysis showed that spacing-aware differences vary with gray-level discretization. Therefore, binWidth and discretization settings should be reported together with voxel-spacing-aware extraction results.

\begin{table}[htbp]
\centering
\small
\caption{Additional software-validation summary.}
\begin{tabular}{p{0.17\linewidth}p{0.22\linewidth}p{0.18\linewidth}p{0.18\linewidth}p{0.17\linewidth}}
\toprule
Experiment & Main comparison & Key metric & Main result & Interpretation \\
\midrule
Resampling baseline comparison & Voxel-spacing-aware extraction vs native and isotropic resampling & Median relative difference; correlation & VS median relative difference vs native 0.307; linear resampling 0.304. & VS behaves as a distinct interpolation-free alternative. \\
Computational profiling & Spacing ratio and ROI size & Runtime vs expansion factor & Observed up to R=7; maximum completed runtime 1.74 s. & Cost scales with finite-volume expansion; warnings are advisable for high R. \\
Rounding sensitivity & Integer vs non-integer spacing ratios & Relative represented-spacing error & Maximum tested relative z-spacing error 0.2. & Non-integer ratios are approximated by integer repeat factors. \\
Binning sensitivity & binWidth 5, 10, 25, 50 & Median relative difference & VS vs native median relative difference range 0.145-0.233. & Discretization interacts with feature-level differences and should be reported. \\
\bottomrule
\end{tabular}
\end{table}


\subsection*{Discussion insert}
These software-validation experiments do not test clinical performance and should not be interpreted as evidence of diagnostic superiority. Instead, they document implementation behavior under controlled phantom-like conditions. Voxel-spacing-aware extraction is best described as an interpolation-free extension of standard PyRadiomics texture computation that preserves the gray levels supplied to texture computation while incorporating physical voxel spacing into selected texture families. Because these outputs extend standard definitions, they should not be reported as IBSI reference values unless and until the definitions are formally standardized.

\subsection*{Limitations insert}
The added experiments are synthetic only and intentionally avoid clinical modelling. They do not establish patient-level robustness or prognostic performance. The finite-volume strategy introduces computational overhead proportional to the expansion factor and approximates non-integer spacing ratios through integer repeat factors. Further validation should include digital phantoms, test-retest data, and independent clinical robustness studies. Application-level clinical behavior is outside the scope of this technical implementation paper and is addressed separately in a companion study the companion clinical study, where the author may insert an approved one-sentence finding here: reserved wording to be inserted by the authors if needed.
